\chapter{GPU-DAMOV: Methodology and Experimental Setup}
\label{ch:method}

This chapter defines the method. Section~\ref{sec:met-pipeline} gives the
pipeline; \S\ref{sec:met-setup} the instruments and workloads;
\S\ref{sec:met-metrics} every metric with its formula;
\S\ref{sec:met-tree} the classifier and taxonomy; \S\ref{sec:met-verdict}
the substrate-verdict and fault mapping; \S\ref{sec:met-model} the placement
model; and \S\ref{sec:met-parity} the component-by-component parity argument
that GPU-DAMOV is complete at DAMOV's granularity (RQ1).

\section{The six-step pipeline}
\label{sec:met-pipeline}

\begin{center}
\code{Capture} $\rightarrow$ \code{Screen} $\rightarrow$ \code{Classify}
$\rightarrow$ \code{Attribute} $\rightarrow$ \code{Validate} $\rightarrow$
\code{Verdict}
\end{center}

\textbf{Capture} runs the workload under the three profilers at locked
clocks: nsys for the timeline, ncu for counters (bounded launch windows,
curated metric set), NVBit for address traces (windowed). Every run also
records whole-run energy (GPU power sampled and integrated), host-side I/O
(storage bytes read, memory-mapped page-ins, peak host RSS), and---always---
trajectory accuracy against ground truth, so that no measurement of a
malfunctioning run can enter the evidence. \textbf{Screen} ranks kernels by
share of GPU time from the timeline; kernels below the attention floor are
screened out, not classified (they receive class G0). \textbf{Classify}
feeds each surviving kernel's metrics to the ordered decision tree
(\S\ref{sec:met-tree}). \textbf{Attribute} joins address traces against the
allocation journal to name each kernel's traffic
(Chapter~\ref{ch:attribution}). \textbf{Validate} stress-tests everything:
threshold sensitivity, unsupervised clustering, cross-device agreement, a
known-truth calibration suite, and profiler-neutrality checks
(Chapters~\ref{ch:validity}--\ref{ch:characterization}). \textbf{Verdict}
maps class $\times$ data-structure persistence $\times$ features to a
substrate, with a named on-GPU fault when the verdict is ``stay.''

\section{Experimental setup}
\label{sec:met-setup}

\subsection{Hardware}

The primary instrument is an NVIDIA RTX~2000 Ada generation GPU
(\code{sm\_89}): 22~SMs, 100\,KB unified L1/shared per SM, 25\,MB L2, 16\,GB
GDDR6. A second, deliberately different microarchitecture---a GeForce MX450
laptop part (\code{sm\_75}, 2\,GB, 25\,W)---serves only for cross-device
validation of the classification (\S\ref{sec:char-crossdev}). The host
workstation runs driver 575.64.05 with CUDA 12.9.1; profilers are Nsight
Compute 2025.2.1.3 and Nsight Systems 2025.3.2; address tracing uses NVBit
1.8~\cite{villa2019nvbit}. This exact combination matters operationally
(binary instrumentation constrains the driver generation, which in turn
constrains profiler versions) and is documented in the artifact so the
environment is reconstructible (Appendix~\ref{app:artifacts}).

\subsection{Locked clocks and measured ceilings}
\label{sec:met-clocks}

GPUs modulate frequency continuously for power and thermals; two runs at
different clocks are not comparable. All measurements lock the graphics
clock at 1620\,MHz and the memory clock at 7001\,MHz. The effect is not
cosmetic: the median run-to-run coefficient of variation (CoV = standard
deviation $\div$ mean over repeated runs) of kernel time drops from
\textbf{49.6\%} unlocked on a thermally limited laptop, and 9.3\% warmed, to
\textbf{0.14\%} locked on the workstation. At 0.14\%, run-to-run noise is
three orders of magnitude below every effect size reported in this thesis.

Roofline ceilings are \emph{measured}, not quoted from datasheets: at the
locked clocks, sustained DRAM bandwidth is $205.0 \pm 0.1$\,GB/s and FP32
throughput is $5445 \pm 3$\,GFLOP/s (microbenchmarks in the artifact). All
roofline placements and DRAM speed-of-light percentages in this thesis are
relative to these measured ceilings.

\subsection{Workloads and datasets}
\label{sec:met-datasets}

The characterization corpus is 27 sequences across four public ground-truth
datasets: KITTI odometry 00--10 (automotive stereo, kilometre
scale)~\cite{geiger2013kitti}; EuRoC MAV MH01--05, V101--103, V201--203
(drone stereo + IMU, room/hall scale)~\cite{burri2016euroc}; TUM RGB-D fr3
sequences including the loop-closure-dense \code{long\_office\_household}
(handheld RGB-D)~\cite{sturm2012tum}; and TUM-VI
(visual-inertial)~\cite{schubert2018tumvi}. The accuracy-validation matrix
(Chapter~\ref{ch:validity}) additionally includes
ICL-NUIM~\cite{handa2014iclnuim}. Sequences span indoor and outdoor,
room-scale and street-scale, loop-closing and deliberately loop-free
trajectories (KITTI 01 and 04 close no loops and serve as a control).
Pipeline modes exercised: odometry-only and full SLAM; synchronous and
asynchronous SLAM; GPU and CPU SLAM back-ends; stereo, RGB-D, monocular, and
inertial sensor configurations. In total the configuration regime
materializes 141 accuracy configurations and 192 profiling-coverage
variants from a small set of human-owned base configurations via a
deterministic mutation pass, so the whole matrix is reproducible
byte-for-byte.

\subsection{Capture discipline}

Three rules are enforced throughout. (1)~\emph{Traces are never used for
timing}: instrumented runs are 5--1000$\times$ slower by design; all timing
comes from uninstrumented locked-clock runs. (2)~\emph{Windows must be
audited}: windowed captures risk missing sparsely-launching kernels, so a
coverage audit diffs every sequence's full launch map against the captured
set and iteratively gap-fills until zero kernels are missing
(\S\ref{sec:attr-coverage}). (3)~\emph{Extensive before intensive}: time,
bytes, and instruction counts are summed across a kernel's launches first
and ratios are formed once; averaging per-launch ratios would weight a
microsecond launch equally with a millisecond one.

\section{Metrics}
\label{sec:met-metrics}

Each metric reduces hardware counters or trace records to one interpretable
feature per kernel.

\subsection{Speed-of-light utilizations}
Nsight Compute reports each subsystem's utilization as a percentage of its
theoretical peak: \emph{memory SoL} (the memory pipeline),
\emph{compute SoL} (the arithmetic pipelines), and \emph{DRAM SoL} (DRAM
bandwidth). DRAM SoL $\geq$ 50\% is read as bandwidth saturation; compute
SoL substantially above memory SoL indicates compute-boundedness.

\subsection{LFMR (GPU form)}
\begin{equation}
\text{LFMR} \;=\; 1 - \text{(L2 hit rate)}.
\label{eq:lfmr}
\end{equation}
DAMOV's CPU definition---misses at the last-level cache divided by misses at
the first---collapses on a two-level hierarchy to the fraction of L2
accesses that proceed to DRAM. LFMR $\approx 1$: the L2 contributes nothing;
the access stream is cache-defeating (PiM-favourable). LFMR $\leq 0.35$:
the L2 is absorbing reuse; near-data placement would forfeit a working
cache.

\subsection{Memory intensity (MPKI)}
\begin{equation}
\text{MPKI} \;=\; \frac{\text{DRAM bytes}/32}
{\text{warp instructions}/1000},
\label{eq:mpki}
\end{equation}
i.e., DRAM 32-byte-sector fetches per thousand \emph{warp} instructions.
(The instruction denominator is warp-level, $32\times$ smaller than
thread-level; conflating the two is a classic factor-of-32 error the
pipeline guards against in tests.)

\subsection{Arithmetic intensity and the roofline}
\label{sec:met-roofline}
\begin{equation}
\text{AI} \;=\; \frac{\text{FLOPs}}{\text{DRAM bytes}},\qquad
\text{FLOPs} = \text{adds} + \text{muls} + 2\times\text{FMAs}.
\label{eq:ai}
\end{equation}
Attainable performance is $\min(P_{\text{peak}},\ \text{AI}\times
B_{\text{peak}})$ with the \emph{measured} ceilings of
\S\ref{sec:met-clocks}~\cite{williams2009roofline}. A kernel under the
sloped segment is memory-bound; under the flat segment, compute-bound. The
FLOP numerator selects the dominant operand type automatically (FP32 for
cuVSLAM; FP16/INT for other adapters' workloads), keeping the roofline
meaningful across workload families.

\subsection{Occupancy and the stall taxonomy}
Occupancy is the achieved fraction of maximum resident warps; below 25\%
the latency-hiding mechanism of \S\ref{sec:bg-gpu} is starved. When a warp
cannot issue, the hardware records \emph{why}; the classifier consumes the
dominant reason: \code{long\_scoreboard} (waiting on DRAM/L2 data---a memory
stall), \code{short\_scoreboard} (shared-memory data), \code{lg/mio/tex
throttle} (memory-pipe queues full), \code{wait} (fixed-latency dependency
chains---\emph{not} memory), \code{barrier} (synchronization). The pair
(dominant stall, occupancy) separates ``memory-latency-bound'' (G4) from
``dependency-bound'' (G7)---two conditions with identical symptoms in a
time-only profile and opposite architectural remedies.

\subsection{Coalescing}
Sectors per request, from counters, with the trace-side check of
sectors-per-warp-access from NVBit (\S\ref{sec:loc-metrics});
$\geq 8$ marks a scattered gather.

\subsection{Locality (trace-side)}
Defined in \S\ref{sec:loc-metrics}: reuse-distance CDFs, footprints, and
inter-launch overlap, computed from global-space addresses only.

\section{The classifier and taxonomy}
\label{sec:met-tree}

\subsection{Thresholds}
The decision tree's thresholds (Table~\ref{tab:thresholds}) are stated once
in the classifier source and imported live into all documentation and the
dashboard, so prose and code cannot drift. They are calibrated on the
characterization corpus and then \emph{frozen}; robustness is established by
the $\pm 25\%$ sensitivity stress (\S\ref{sec:met-sensitivity}) and the
known-truth calibration suite (\S\ref{sec:char-calibration}) rather than by
tuning.

\begin{table}[t]
\centering
\caption{Classifier thresholds (values as frozen in
\code{analysis/classify.py}).}
\label{tab:thresholds}
\begin{tabular}{@{}llp{0.52\textwidth}@{}}
\toprule
threshold & value & meaning \\
\midrule
\code{sol\_hi} & 40\,\% & a speed-of-light value counted as ``high'' \\
\code{sol\_ratio} & 1.5$\times$ & margin by which compute SoL must exceed
memory SoL for compute-boundedness \\
\code{dram\_sat} & 50\,\% & DRAM SoL treated as bandwidth saturation \\
\code{lfmr\_hi} / \code{lfmr\_lo} & 0.40 / 0.35 & L2 ``not helping'' /
``earning its keep'' \\
\code{sect\_scatter} & 8 & sectors/request marking a scattered gather \\
\code{occ\_low} / \code{occ\_low\_dep} & 25\,\% / 30\,\% & occupancy below
which latency (G4) / dependency (G7) stalls cannot be hidden \\
\bottomrule
\end{tabular}
\end{table}

\subsection{The ordered rules}
First matching rule assigns the class:
\begin{enumerate}
\item \textbf{G5 compute-bound}: compute SoL $\geq$ 40\% and $\geq$
1.5$\times$ memory SoL.
\item \textbf{G6 on-chip-bound}: a shared-memory/MIO stall dominates while
DRAM is not saturated.
\item \textbf{G1 bandwidth-bound}: DRAM SoL $\geq$ 50\%.
\item \textbf{G2 coalescing-bound}: memory-limited and sectors/request
$\geq$ 8.
\item \textbf{G3 L2-reuse-bound}: memory-limited and LFMR $<$ 0.35.
\item \textbf{G4 latency-bound}: a memory stall dominates at occupancy
$<$ 25\% with DRAM unsaturated.
\item \textbf{G7 dependency-bound}: a dependency stall (\code{wait},
\code{barrier}, \ldots) dominates at low occupancy with neither memory nor
compute high.
\item \textbf{G0 no-signal}: nothing dominant (tiny or launch-tax kernels;
screened rather than architecturally interpreted).
\end{enumerate}

Class G7 was not part of the initial hypothesis; it emerged from cuVSLAM
kernels that stall on instruction dependencies at low occupancy with memory
idle. That the taxonomy grew a class under pressure from data is by design
---DAMOV's own classes were likewise derived, not decreed.

\begin{table}[t]
\centering
\caption{The GPU bottleneck taxonomy and its default near-data affinity.}
\label{tab:taxonomy}
\begin{tabular}{@{}llp{0.35\textwidth}@{}}
\toprule
class & meaning & PiM affinity \\
\midrule
G0 & no dominant signal / sub-floor & --- (candidate for CPU) \\
G1 & DRAM bandwidth saturated & \textbf{strong} \\
G2 & scattered access wastes bandwidth & \textbf{conditional} (scatter PiM
or layout fix) \\
G3 & L2 reuse effective & weak (keep the cache) \\
G4 & memory-latency stalls, low occupancy & strong if cache-defeating \\
G5 & compute-bound & none (keep on GPU) \\
G6 & on-chip/shared-memory-bound & none \\
G7 & dependency stalls, memory idle & none (fix occupancy first) \\
\bottomrule
\end{tabular}
\end{table}

\subsection{Sensitivity analysis}
\label{sec:met-sensitivity}
Every kernel is re-classified with all thresholds scaled by $\pm 25\%$. A
kernel whose class changes under this perturbation is flagged
\emph{borderline} and cannot carry high confidence into the verdict stage;
every headline kernel in this thesis is threshold-stable. The observed
flips concentrate exactly where physics predicts: at the L2-capacity
crossover, where a working set close to 25\,MB moves between G1 and G3.

\section{From class to verdict: substrate and fault mapping}
\label{sec:met-verdict}

Class alone does not determine placement; the join with data-structure
\emph{persistence} (Chapter~\ref{ch:attribution}) does. Persistence takes
three values, measured from the allocation journal and stage map:
\emph{streaming} (per-frame data, consumed once), \emph{hot-persistent}
(resident state reused every frame, e.g.\ the BA linear system), and
\emph{cold-persistent} (session-lifetime state touched episodically, e.g.\
the keyframe database). Table~\ref{tab:verdict} gives the mapping.

\begin{table}[t]
\centering
\caption{Substrate verdict mapping.}
\label{tab:verdict}
\begin{tabular}{@{}lll@{}}
\toprule
condition & affinity & substrate \\
\midrule
G1 $\wedge$ streaming & strong & near-sensor SRAM (consume before DRAM) \\
G1 $\wedge$ otherwise & strong & DRAM-PiM (near-bank) \\
G2 & conditional & scatter-capable PiM, or data-layout fix first \\
G4 $\wedge$ cache-defeating & strong & near-memory compute \\
cold-persistent $\wedge$ large & strong & in/near-storage (ISP) \\
G3 & weak & GPU (a larger cache wins) \\
G5 / G6 / G7 & none & GPU \\
tiny (occ $<$ 8\%, $<$ 1\,ms) & --- & CPU/host \\
\bottomrule
\end{tabular}
\end{table}

When the verdict is ``stay on GPU,'' the same evidence names the fault:
spill-dominated traffic $\Rightarrow$ register-pressure/compiler fault
(the traffic is scratch, not data); high sectors/request $\Rightarrow$
data-layout fault; dependency stalls at low occupancy $\Rightarrow$
launch-configuration fault; a flat reuse CDF $\Rightarrow$ structurally
cache-immune (no cache-size remedy exists). This dual output---substrate
candidacy \emph{or} a named fix---is what makes the characterization
actionable in both directions.

\section{The placement model}
\label{sec:met-model}

To bound the payoff of offloading without a simulator, a first-order
analytical model: for a kernel of GPU time $t$ with memory-bound fraction
$m$, executed on a near-data substrate with internal-bandwidth advantage $k$
and relative compute throughput $c$,
\begin{equation}
t_{\text{pim}} \;=\; t\,\frac{1-m}{c} \;+\; t\,\frac{m}{k}.
\label{eq:placement}
\end{equation}
A kernel is offloaded only if its affinity is admitted by the scenario and
$t_{\text{pim}} < t$. Two scenarios bound the claim: \emph{conservative}
($k{=}4$, $c{=}0.5$, strong-affinity only) and \emph{moderate} ($k{=}8$,
$c{=}0.75$, strong + conditional). An energy variant applies the same split
to the measured energy baseline. A standing rule constrains all of this:
model-derived and (future) simulator-derived numbers are reported only as
\emph{deltas against the measured GPU baseline}, never as absolute
performance claims.

\section{DAMOV parity: the completeness argument}
\label{sec:met-parity}

RQ1 asks whether every DAMOV component has a GPU analogue here. The
correspondence, element by element:

\begin{itemize}
\item \emph{Step 1, screening.} DAMOV: VTune top-down, keep functions with
memory-bound share $>$30\% and $\geq$3\% of cycles. Here: nsys time-share
plus the ncu SoL pair; sub-floor kernels are screened (G0), not classified.
\item \emph{Step 2, locality.} DAMOV: architecture-independent spatial and
temporal locality from CPU address streams, clustered. Here: NVBit per-warp
traces $\rightarrow$ reuse-distance CDFs (temporal) and sectors-per-warp
coalescing (the GPU's spatial analogue), \emph{memory-space filtered} so
compiler scratch cannot pollute locality (\S\ref{sec:loc-correction}).
\item \emph{Step 3, intervention.} DAMOV: simulated host / host+prefetch /
NDP across a 1--256-core sweep, with classes defined by the response. Here,
two measured axes replace the unavailable core sweep (a GPU grid is compiled
into the launch; SM count is not variable on this part): (a) a
\emph{clock-domain sweep}---core and memory clocks scaled independently
(1620$\rightarrow$810\,MHz core = 0.5$\times$ compute; 7001$\rightarrow$5001
= 0.71$\times$ DRAM)---yields falsifiable per-class predictions (G1/G2 track
the memory clock, G3/G5/G6/G7 the core clock, G4 neither strongly); (b) the
\emph{measured} reuse-distance hit-CDF answers the cache-capacity question
directly, per kernel, from 64\,KiB to 48\,MiB, where DAMOV inferred it from
aggregate sweeps. The third, simulated axis (NDP configurations for
Accel-Sim) is generated from the verdicts and deferred to the follow-on
study.
\item \emph{Classification.} DAMOV: six classes from a threshold tree over
temporal locality, LFMR trend, MPKI, AI. Here: eight classes (G0--G7) from
the ordered tree of \S\ref{sec:met-tree}; the deviations (a coalescing
class; a G4/G7 occupancy split; no L3-contention class because the L2 is the
last level) are each forced by a physical difference of the GPU.
\item \emph{Robustness.} DAMOV: frozen thresholds validated on 100 held-out
functions (97\% correct); cache-size sweeps; an independent clustering
algorithm. Here: the $\pm$25\% sensitivity stress; a known-truth
\emph{calibration suite} of eight archetype kernels written to embody each
class and classified blind (\S\ref{sec:char-calibration}); cross-device
agreement across two microarchitectures (\S\ref{sec:char-crossdev}); and
\emph{two} independent clustering algorithms (k-means and Ward hierarchical)
recovering the taxonomy (\S\ref{sec:char-clustering}).
\end{itemize}

Where DAMOV's breadth was 144 applications, this study is deliberately
deep rather than wide---one production application, 49 kernels, 27
sequences, 192 configuration variants---with breadth provided
architecturally: the harness's adapter interface runs arbitrary GPU
workloads through the identical pipeline, and the calibration archetypes
serve as the suite-artifact analogue.
